\documentclass[a4paper,12pt]{article}

\usepackage[T2A]{fontenc}
\usepackage[utf8]{inputenc}
\usepackage[russian]{babel}

\usepackage{amsmath,amssymb,mathtools}
\usepackage{helvet}
\renewcommand{\familydefault}{\sfdefault}

\usepackage{icomma}
\usepackage[a4paper,margin=2.2cm,headheight=16pt]{geometry}
\usepackage{microtype}
\usepackage{tikz}
\usepackage{tabularx,booktabs}
\usepackage{enumitem}
\usepackage{hyperref}
\usepackage{parskip}
\usepackage{fancyhdr}
\usepackage{setspace}
\usepackage{xcolor}

\definecolor{accent}{HTML}{008080}
\definecolor{darkaccent}{HTML}{004D4D}
\definecolor{textgray}{HTML}{333333}
\definecolor{softgray}{HTML}{EEF3F3}
\definecolor{linegray}{HTML}{B8C8C8}

\hypersetup{
    colorlinks=true,
    linkcolor=darkaccent,
    urlcolor=darkaccent,
    pdftitle={Чек-лист. Повторение вычислений},
    pdfauthor={Лёвин Артём Александрович}
}

\setstretch{1.35}
\setlength{\parindent}{0pt}
\setlength{\parskip}{0.55em}

\setlist[itemize]{
    leftmargin=1.5em,
    itemsep=0.25em,
    topsep=0.25em
}
\setlist[enumerate]{
    leftmargin=1.8em,
    itemsep=0.45em,
    topsep=0.35em
}

\pagestyle{fancy}
\fancyhf{}
\fancyhead[L]{\small\color{textgray}Повторение вычислений}
\fancyhead[R]{\small\color{textgray}7 класс}
\fancyfoot[C]{\small\color{textgray}\thepage}
\renewcommand{\headrulewidth}{0.4pt}
\renewcommand{\headrule}{\hbox to\headwidth{\color{linegray}\leaders\hrule height \headrulewidth\hfill}}

\newcommand{\sectionline}{%
    \vspace{-0.25em}
    {\color{accent}\rule{\linewidth}{0.7pt}}
    \vspace{0.3em}
}

\newcommand{\checkitem}{\item[\textcolor{accent}{$\checkmark$}]}

\newcommand{\important}[1]{\textcolor{darkaccent}{\textbf{#1}}}

\begin{document}

% ============================================================
% ТИТУЛЬНЫЙ ЛИСТ
% ============================================================

\begin{titlepage}
    \thispagestyle{empty}

    \begin{center}
        \vspace*{2.7cm}

        {\color{darkaccent}\Large\bfseries ЧЕК-ЛИСТ}

        \vspace{0.8cm}

        {\Huge\bfseries Повторение вычислений}

        \vspace{0.45cm}

        {\Large
        Десятичные и обыкновенные дроби\\
        Отрицательные числа. Проценты}

        \vspace{1.5cm}

        {\color{accent}\rule{0.72\linewidth}{1pt}}

        \vspace{1.5cm}

        {\large
        Лёвин Артём Александрович\\[0.35em]
        эксклюзивно для Кирилла Зиновьева}

        \vspace{0.9cm}

        {\large 7 класс}

        \vfill

        {\large \today}

        \vspace{2.4cm}
    \end{center}

    \begin{tikzpicture}[remember picture,overlay]
        \draw[
            color=accent!55,
            line width=1.2pt
        ]
        ([xshift=1.6cm,yshift=1.55cm]current page.south west)
        .. controls
        ([xshift=5.0cm,yshift=2.15cm]current page.south west)
        and
        ([xshift=-5.5cm,yshift=0.95cm]current page.south east)
        ..
        ([xshift=-1.6cm,yshift=1.6cm]current page.south east);
    \end{tikzpicture}
\end{titlepage}

% ============================================================
% ОСНОВНОЙ ЧЕК-ЛИСТ
% ============================================================

\section*{\color{darkaccent}1. Численные выражения}
\sectionline

\textbf{Численное выражение} --- запись, составленная из чисел, знаков действий и, при необходимости, скобок.

\textbf{Значение численного выражения} --- число, которое получается после выполнения всех действий.

\[
6{,}965+23{,}3=30{,}265.
\]

\begin{itemize}
    \checkitem Перед вычислением определите порядок действий.
    \checkitem Выполняйте вычисления аккуратно и проверяйте положение запятой.
    \checkitem После получения результата оцените, похож ли он по величине на ожидаемый.
\end{itemize}

\subsection*{\color{accent}Порядок действий}

Если скобок нет:
\[
\boxed{\text{умножение и деление}}
\quad\longrightarrow\quad
\boxed{\text{сложение и вычитание}}.
\]

Действия одной ступени выполняются слева направо.

Если есть скобки:
\[
\boxed{\text{скобки}}
\quad\longrightarrow\quad
\boxed{\text{умножение и деление}}
\quad\longrightarrow\quad
\boxed{\text{сложение и вычитание}}.
\]

Например:
\[
8+12:3\cdot2=8+4\cdot2=16.
\]

\[
(8+12):4=20:4=5.
\]

% ============================================================

\section*{\color{darkaccent}2. Десятичные дроби}
\sectionline

\subsection*{\color{accent}Сложение и вычитание}

При записи столбиком десятичные дроби располагают так, чтобы
\important{запятые находились друг под другом}.

Недостающие разряды можно дополнить нулями:
\[
50{,}4=50{,}40.
\]

Пример:
\[
50{,}40-6{,}98=43{,}42.
\]

\begin{itemize}
    \checkitem Запишите числа запятая под запятой.
    \checkitem При необходимости допишите нули справа.
    \checkitem Выполните обычное сложение или вычитание по разрядам.
    \checkitem Запятую результата поставьте под запятыми исходных чисел.
\end{itemize}

\subsection*{\color{accent}Умножение}

При умножении десятичных дробей числа удобно сначала перемножить как целые.

После вычисления произведения нужно подсчитать
\important{общее количество цифр после запятых} в обоих множителях.

Пример:
\[
6{,}5\cdot1{,}22.
\]

У первого множителя одна цифра после запятой, у второго --- две. Всего:
\[
1+2=3.
\]

Как целые числа:
\[
65\cdot122=7930.
\]

Отделяем справа три цифры:
\[
6{,}5\cdot1{,}22=7{,}930=7{,}93.
\]

\subsection*{\color{accent}Деление}

При делении столбиком удобно добиться того, чтобы
\important{делитель был целым числом}.

Если делитель содержит запятую, запятую сдвигают вправо одновременно
в делимом и делителе на одинаковое количество разрядов.

Например:
\[
16{,}94:2{,}8
=
169{,}4:28.
\]

Это преобразование допустимо, поскольку:
\[
\frac{16{,}94}{2{,}8}
=
\frac{16{,}94\cdot10}{2{,}8\cdot10}.
\]

Если в делителе две цифры после запятой, потребуется умножение на \(100\):
\[
16{,}94:0{,}28
=
1694:28.
\]

\important{Правило.}
Количество разрядов, на которое сдвигается запятая, определяется количеством
цифр после запятой в делителе.

Пример обычного деления:
\[
53{,}4:15=3{,}56.
\]

% ============================================================

\section*{\color{darkaccent}3. Обыкновенные дроби}
\sectionline

Для дроби
\[
\frac{a}{b}
\]
число \(a\) называется \textbf{числителем}, число \(b\) ---
\textbf{знаменателем}.

Всегда выполняется ограничение:
\[
\boxed{b\ne0}.
\]

\subsection*{\color{accent}Сокращение дробей}

Числитель и знаменатель можно разделить на одно и то же ненулевое число.

Например:
\[
\frac{26}{24}
=
\frac{13}{12}.
\]

Перед окончательной записью ответа полезно проверить:

\begin{itemize}
    \checkitem можно ли сократить дробь;
    \checkitem получилась ли неправильная дробь;
    \checkitem требуется ли выделить целую часть.
\end{itemize}

% ============================================================

\section*{\color{darkaccent}4. Сложение и вычитание дробей}
\sectionline

\subsection*{\color{accent}Способ 1. Общий знаменатель}

Для сложения:
\[
\frac{5}{6}+\frac14.
\]

Наименьший общий знаменатель чисел \(6\) и \(4\):
\[
12.
\]

Тогда:
\[
\frac56=\frac{10}{12},
\qquad
\frac14=\frac3{12}.
\]

Следовательно:
\[
\frac56+\frac14
=
\frac{10}{12}+\frac3{12}
=
\frac{13}{12}
=
1\frac1{12}.
\]

\important{Знаменатели при сложении напрямую не складываются.}

\subsection*{\color{accent}Способ 2. Перекрёстное умножение}

Общая формула:
\[
\frac ab+\frac cd
=
\frac{ad+bc}{bd},
\qquad b\ne0,\quad d\ne0.
\]

Для вычитания:
\[
\frac ab-\frac cd
=
\frac{ad-bc}{bd}.
\]

Этот способ особенно удобен, когда знаменатели небольшие или содержат буквы.

% ============================================================

\section*{\color{darkaccent}5. Смешанные и неправильные дроби}
\sectionline

\textbf{Неправильная дробь} --- дробь, у которой числитель больше знаменателя
или равен ему.

Например:
\[
\frac{13}{12}=1\frac1{12}.
\]

\subsection*{\color{accent}Из смешанной дроби в неправильную}

\[
3\frac27
=
\frac{3\cdot7+2}{7}
=
\frac{23}{7}.
\]

Общий алгоритм:
\[
a\frac bc
=
\frac{ac+b}{c}.
\]

\subsection*{\color{accent}Целое число как дробь}

Любое целое число можно записать со знаменателем \(1\):
\[
5=\frac51,
\qquad
10=\frac{10}{1}.
\]

При необходимости целое число можно представить и с другим знаменателем:
\[
5=\frac{35}{7}.
\]

Это удобно, например, при вычислении:
\[
5-3\frac27.
\]

Преобразуем:
\[
5=\frac{35}{7},
\qquad
3\frac27=\frac{23}{7}.
\]

Тогда:
\[
5-3\frac27
=
\frac{35}{7}-\frac{23}{7}
=
\frac{12}{7}
=
1\frac57.
\]

% ============================================================

\section*{\color{darkaccent}6. Умножение дробей}
\sectionline

Для умножения дробей:
\[
\frac ab\cdot\frac cd
=
\frac{ac}{bd},
\qquad b\ne0,\quad d\ne0.
\]

Пример:
\[
\frac49\cdot\frac38.
\]

Можно сразу сократить множители:
\[
\frac49\cdot\frac38
=
\frac{1}{3}\cdot\frac{1}{2}
=
\frac16.
\]

\important{Полезный приём:}
если сокращение видно заранее, выполняйте его до перемножения.
Числа останутся меньше, поэтому риск арифметической ошибки снизится.

Сокращение выполняют между множителями, расположенными по разные стороны
дробной черты:
\[
\frac{\cancel{4}}{9}\cdot\frac{3}{\cancel{8}}
\]
в таком оформлении потребовался бы дополнительный пакет, поэтому в рабочей
записи достаточно явно указать деление соответствующих множителей.

Например:
\[
\frac49\cdot\frac38
=
\frac{1}{3}\cdot\frac12.
\]

% ============================================================

\section*{\color{darkaccent}7. Деление дробей}
\sectionline

При делении на дробь деление заменяется умножением, а дробь после знака
деления заменяется обратной:

\[
\frac ab:\frac cd
=
\frac ab\cdot\frac dc,
\]
где
\[
b\ne0,\qquad c\ne0,\qquad d\ne0.
\]

Условие \(c\ne0\) особенно важно: делить на ноль нельзя.

Пример:
\[
\frac58:\frac9{10}
=
\frac58\cdot\frac{10}{9}
=
\frac{25}{36}.
\]

\subsection*{\color{accent}Деление смешанных дробей}

Перед умножением и делением смешанные дроби удобно перевести в неправильные.

Например:
\[
2\frac67:1\frac37.
\]

Преобразуем:
\[
2\frac67=\frac{20}{7},
\qquad
1\frac37=\frac{10}{7}.
\]

Тогда:
\[
\frac{20}{7}:\frac{10}{7}
=
\frac{20}{7}\cdot\frac7{10}
=
2.
\]

\important{Алгоритм:}
\begin{enumerate}
    \item Переведите смешанные дроби в неправильные.
    \item Замените деление умножением.
    \item Переверните дробь, стоящую после знака деления.
    \item Сократите множители.
    \item Перемножьте оставшиеся числители и знаменатели.
    \item При необходимости выделите целую часть.
\end{enumerate}

% ============================================================

\section*{\color{darkaccent}8. Отрицательные числа}
\sectionline

\subsection*{\color{accent}Сложение и вычитание}

Примеры:
\[
10+2=12,
\]
\[
10-2=8,
\]
\[
2-10=-8,
\]
\[
-2-10=-12,
\]
\[
-2+10=8.
\]

Если складываются числа разных знаков, удобно:

\begin{enumerate}
    \item сравнить их модули;
    \item из большего модуля вычесть меньший;
    \item поставить знак числа с большим модулем.
\end{enumerate}

Например:
\[
-13+8=-5.
\]

\subsection*{\color{accent}Умножение и деление. Правило знаков}

\[
(+)\cdot(+)=+,
\qquad
(+)\cdot(-)=-,
\]
\[
(-)\cdot(+)=-,
\qquad
(-)\cdot(-)=+.
\]

Те же правила действуют при делении.

Например:
\[
10\cdot(-2)=-20,
\]
\[
(-10)\cdot(-2)=20.
\]

Если отрицательное число является множителем, его удобно заключать в скобки:
\[
10\cdot(-2).
\]

% ============================================================

\section*{\color{darkaccent}9. Проценты}
\sectionline

\textbf{Один процент} --- одна сотая часть величины:
\[
1\%=\frac1{100}=0{,}01.
\]

Чтобы перевести \(p\%\) в десятичную дробь:
\[
\boxed{p\%=\frac{p}{100}}.
\]

Примеры:
\[
40\%=0{,}4,
\]
\[
138\%=1{,}38,
\]
\[
0{,}43\%=0{,}0043.
\]

Обратите внимание:
\[
120\%=1{,}2.
\]

Значит, процент может соответствовать числу, большему единицы.

\subsection*{\color{accent}Как найти процент от числа}

Чтобы найти \(p\%\) от числа \(A\), удобно выполнить два шага:

\[
p\%
\longrightarrow
\frac{p}{100}
\longrightarrow
A\cdot\frac{p}{100}.
\]

То есть:
\[
\boxed{p\%\text{ от }A=A\cdot\frac{p}{100}}.
\]

Пример:
\[
40\%\text{ от }15
=
15\cdot0{,}4
=
6.
\]

Ещё один пример:
\[
120\%\text{ от }8
=
8\cdot1{,}2
=
9{,}6.
\]

% ============================================================

\section*{\color{darkaccent}10. Умножение и деление на \(10\), \(100\), \(1000\)}
\sectionline

При умножении на степень числа \(10\) запятая перемещается вправо:

\[
4{,}27\cdot10=42{,}7,
\]
\[
4{,}27\cdot100=427,
\]
\[
4{,}27\cdot1000=4270.
\]

При делении запятая перемещается влево:

\[
138:100=1{,}38,
\]
\[
13{,}8:1000=0{,}0138.
\]

Если для записи результата разрядов недостаточно, добавляются нули.

% ============================================================

\section*{\color{darkaccent}11. Финальная самопроверка вычислений}
\sectionline

Перед записью ответа проверьте:

\begin{itemize}
    \checkitem Верно ли определён порядок действий?
    \checkitem Совпадают ли разряды при сложении и вычитании десятичных дробей?
    \checkitem Правильно ли поставлена запятая после умножения?
    \checkitem Стал ли делитель целым при делении десятичных дробей?
    \checkitem Есть ли возможность сократить обыкновенную дробь?
    \checkitem Переведены ли смешанные дроби в неправильные перед умножением или делением?
    \checkitem Перевёрнута ли дробь после знака деления?
    \checkitem Учтено ли правило знаков?
    \checkitem Проценты разделены на \(100\)?
    \checkitem Исключено ли деление на ноль?
    \checkitem Разумна ли величина полученного ответа?
\end{itemize}

% ============================================================
% ТРЕНИРОВОЧНЫЙ БЛОК
% ============================================================

\clearpage
\section*{\color{darkaccent}Тренировочный блок}
\sectionline

\textbf{Уровень:} повторение базовых вычислительных навыков с постепенным
усложнением.

\subsection*{\color{accent}Средний уровень}

\begin{enumerate}
    \item Вычислите:
    \[
    34{,}72+8{,}905-16{,}4.
    \]

    \item Вычислите:
    \[
    4{,}8\cdot1{,}25.
    \]

    \item Найдите:
    \[
    35\%\text{ от }240.
    \]
\end{enumerate}

\subsection*{\color{accent}Выше среднего}

\begin{enumerate}[resume]
    \item Вычислите:
    \[
    18{,}36:2{,}4.
    \]

    \item Вычислите и сократите результат:
    \[
    \frac{7}{12}+\frac{5}{18}.
    \]

    \item Вычислите:
    \[
    6-2\frac58.
    \]

    \item Вычислите:
    \[
    \frac{14}{15}\cdot\frac{25}{28}.
    \]

    \item Вычислите:
    \[
    3\frac38:1\frac7{20}.
    \]

    \item Найдите значение выражения:
    \[
    -18+7\cdot(-3)-(-25).
    \]
\end{enumerate}

\subsection*{\color{accent}Сложный уровень}

\begin{enumerate}[resume]
    \item Вычислите:
    \[
    \left(
        2\frac13-\frac56
    \right)
    :
    \frac34
    +
    12{,}5\%\text{ от }16.
    \]
\end{enumerate}

% ============================================================
% ОТВЕТЫ
% ============================================================

\clearpage
\section*{\color{darkaccent}Ответы к тренировочному блоку}
\sectionline

\renewcommand{\arraystretch}{1.35}

\begin{table}[ht]
    \centering
    \caption{Ответы к упражнениям}
    \begin{tabularx}{\linewidth}{@{}>{\bfseries}c X@{}}
        \toprule
        № & Ответ \\
        \midrule
        1 & \(27{,}225\) \\
        2 & \(6\) \\
        3 & \(84\) \\
        4 & \(7{,}65\) \\
        5 & \(\dfrac{31}{36}\) \\
        6 & \(3\dfrac38\) \\
        7 & \(\dfrac56\) \\
        8 & \(2\dfrac12\) \\
        9 & \(-14\) \\
        10 & \(4\) \\
        \bottomrule
    \end{tabularx}
\end{table}

\vfill

\begin{center}
    {\color{accent}\rule{0.55\linewidth}{0.7pt}}

    \vspace{0.6em}

    {\small\color{textgray}
    Сначала проверяйте алгоритм, затем арифметику и только после этого
    записывайте окончательный ответ.}
\end{center}

\end{document}